\section{Thảo luận}\label{sec:discussion}

\paragraph{Chữ ký test nói được gì và không nói được gì với giảng viên.}
Quy trình kinh điển---OAV theo test case, phân cụm K-means hoặc Hamming, luật NẾU--THÌ quy nạp---nhận ra tốt lỗi phổ biến nhất của khóa học này: lỗi văn bản in tạo chữ ký đặc trưng, và trên bản sửa thật, cụm theo chữ ký đạt ARI trung bình {{real_ari_outcomes}}. Nhưng biểu diễn này không tách được các cơ chế làm trượt cùng những test. Trên các mutant có kiểm soát của cùng các chương trình, nơi mọi cơ chế đều có mặt, cụm của nó hầu như không khớp với cơ chế (ARI {{injected_ari_outcomes}}), và ngay cả một bộ phân loại hoàn hảo chỉ đọc chữ ký cũng chỉ gán đúng tối đa {{injected_H1a_pct}}\% item. Một dashboard dựng trên chữ ký test vì vậy báo cáo chính xác các lỗi phổ biến, thường ở mức đặc tả, và gộp lẫn các lỗi khái niệm ít gặp hơn---điều kiện rẽ nhánh sai, thiếu bước, định dạng in---vốn chính là lý do để phân cụm ngay từ đầu.

\paragraph{Bằng chứng output là cải tiến rẻ nhất, không phải mô hình mới.}
Mọi thứ giúp khôi phục cơ chế tốt hơn đều đã có sẵn trong output của hệ thống chấm: quan hệ giữa output thực và output mong đợi, và trong OAV độ lệch là kiểu sai lệch số và chữ. Trên dữ liệu kiểm soát, điều này làm ARI tăng hơn ba lần ở cùng ngân sách phân cụm, và trên bản sửa thật nó nâng mức khôi phục theo loại của các cơ chế thiểu số. Các mô tả phong phú hơn của chúng tôi không vượt quan hệ output đơn giản khi phân cụm (H3), nên chúng tôi không khẳng định có phương pháp phân cụm tốt hơn. Giá trị của chúng nằm ở chỗ khác: chúng không phụ thuộc bài, nên luật học trên một số bài tập áp dụng được cho bài khác. Luật trên các mô tả độ lệch gộp giữ macro-F1 {{injected_unseen_problems_evidence_ila2_macro_f1}} trên bài tập mới của benchmark kiểm soát, trong khi luật trên tóm tắt kết quả gần như bằng 0; còn chữ ký theo từng test thậm chí không thể so sánh giữa các bài tập.

\paragraph{Luật như những giả thuyết đã kiểm chứng.}
ILA và ILA-2 tạo ra đúng loại luật sản xuất mà cách phát biểu kinh điển yêu cầu, và hệ số phạt của ILA-2 đổi một chút độ chính xác lấy độ phủ lớn hơn nhiều. Chúng tôi trình bày kết luận của luật như một giả thuyết cơ chế kèm độ chính xác đo trên dữ liệu giữ lại, và cho hệ thống từ chối trả lời khi không luật nào khớp. Trong ứng dụng học tập của chúng tôi, mỗi cụm của một lớp giờ hiển thị giả thuyết như vậy, tỷ lệ thành viên khớp luật, một phép kiểm tra một phút để xác nhận hoặc bác bỏ với sinh viên, và một gợi ý giảng lại. Tuy nhiên, trên bản sửa thật, ưu thế của luật bằng chứng so với luật tóm tắt kết quả chưa đáng tin cậy về thống kê với {{real_seen_problems_evidence_ila2_n}} item kiểm tra niêm phong, nên các luật này nên được xem là điểm khởi đầu cho phán đoán của giảng viên.

\paragraph{Embedding tìm ra tác giả trước khi tìm ra cơ chế.}
Một embedding mã đa dụng mạnh đã gom các mutant theo chương trình gốc của chúng: những thay đổi một token quyết định hành vi hầu như không dịch chuyển vector, trong khi tên biến, cách trình bày và chiến lược lời giải chi phối nó. Điều này không mâu thuẫn với các thành công trước đây của embedding được huấn luyện cho đúng nhiệm vụ \citep{shi2021misconception,piech2015embeddings}; nó cho thấy một embedding đóng băng cần bằng chứng thực thi hoặc được huấn luyện nhạy với cơ chế trước khi cụm của nó có thể được đọc như các nhóm lỗi.

\paragraph{Đánh giá dựa trên bản sửa.}
Bản sửa tối thiểu đã kiểm chứng biến {{n_pairs}} sự kiện sửa tự nhiên thành các nhãn có thể audit, với chi phí {{probes_total}} lần thực thi chương trình và không tốn giờ gán nhãn thủ công. Hai phát hiện có ý nghĩa ngoài phạm vi nghiên cứu này. Thứ nhất, {{multi_share}}\% bản sửa tối thiểu thay đổi nhiều hơn một loại mã, nên việc chia các bài sai thành các cụm đơn cơ chế tự thân đã là một sự giản lược; cách gom chồng lấn hoặc đa nhãn cần được quan tâm hơn. Thứ hai, một benchmark kiểm soát xây theo cùng taxonomy làm lộ các giới hạn mà phân bố thật mất cân bằng che giấu, và hai benchmark bất đồng đúng ở chỗ loại chiếm đa số của dữ liệu thật vốn dễ nhận ra.

%%INCLUDE sections/limitations.tex
